% Self-contained LaTeX source. No external figures or bibliography files.
% Compile with:
%   pdflatex bounded_domain_sat_and_amortized_reasoning.tex
%   pdflatex bounded_domain_sat_and_amortized_reasoning.tex
% Alternatively: latexmk -pdf bounded_domain_sat_and_amortized_reasoning.tex
% The author field is intentionally left blank for the recipient to edit.
\documentclass[11pt,a4paper]{article}
\usepackage[T1]{fontenc}
\usepackage[utf8]{inputenc}
\usepackage[margin=25mm,headheight=14pt]{geometry}
\usepackage{amsmath,amssymb,amsthm,mathtools}
\usepackage{newtxtext,newtxmath}
\usepackage{microtype}
\usepackage{graphicx}
\usepackage{xcolor}
\usepackage{tikz}
\usetikzlibrary{arrows.meta,positioning,calc,fit,backgrounds}
\usepackage{caption}
\usepackage{placeins}
\usepackage[authoryear,round]{natbib}
\usepackage{fancyhdr}
\usepackage{hyperref}
\hypersetup{hidelinks,pdftitle={Bounded-Domain SAT, Compiled Knowledge, and Amortized Reasoning},pdfsubject={Finite-domain compilation, neural computation, oracle complexity, and System 1/System 2}}
\urlstyle{same}
\pagestyle{fancy}
\fancyhf{}
\fancyhead[L]{\small Bounded-domain SAT and amortized reasoning}
\fancyhead[R]{\small Perspective}
\fancyfoot[C]{\thepage}
\renewcommand{\headrulewidth}{0.3pt}
\setlength{\parindent}{1.15em}
\setlength{\parskip}{0.2em}
\setlength{\bibsep}{0.45em}
\captionsetup{font=small,labelfont=bf,skip=9pt}
\emergencystretch=2em
\widowpenalty=10000
\clubpenalty=10000
\allowdisplaybreaks[1]
\numberwithin{equation}{section}

\newcommand{\Pclass}{\mathsf{P}}
\newcommand{\NPclass}{\mathsf{NP}}
\newcommand{\coNP}{\mathsf{coNP}}
\newcommand{\Ppoly}{\mathsf{P/poly}}
\newcommand{\SAT}{\mathrm{SAT}}
\newcommand{\bits}{\{0,1\}}
\newcommand{\ind}{\mathbf{1}}
\newcommand{\poly}{\operatorname{poly}}
\newcommand{\ReLU}{\operatorname{ReLU}}
\newcommand{\val}{\operatorname{val}}
\newcommand{\Compile}{\operatorname{Compile}}
\newcommand{\Answer}{\operatorname{Answer}}
\newcommand{\restr}{\mathbin{\upharpoonright}}
\newcommand{\E}{\mathbb{E}}
\DeclareMathOperator{\Prb}{Pr}

\theoremstyle{plain}
\newtheorem{lemma}{Lemma}[section]
\newtheorem{proposition}[lemma]{Proposition}
\newtheorem{corollary}[lemma]{Corollary}
\theoremstyle{definition}
\newtheorem{definition}[lemma]{Definition}
\theoremstyle{remark}
\newtheorem{remark}[lemma]{Remark}

\tikzset{
  block/.style={draw=black!65,rounded corners=2pt,fill=black!3,
    align=center,inner sep=7pt,font=\small,minimum height=12mm},
  memory/.style={block,fill=black!9},
  flow/.style={-{Latex[length=2.0mm]},semithick,draw=black!75},
  optional/.style={flow,dashed},
  note/.style={font=\footnotesize,align=center,inner sep=2pt},
  band/.style={font=\small\bfseries,anchor=west}
}

\title{\vspace{-1.0cm}\textbf{Bounded-Domain SAT, Compiled Knowledge,\\and Amortized Reasoning}\\[0.5em]
\large A perspective on resource accounting for learned reasoning}
\author{}
\date{7 November 2026}

\begin{document}
\maketitle
\vspace{-0.7cm}

\begin{abstract}
Learned reasoning redistributes computation between acquiring knowledge and using it. This perspective brings together established ideas from knowledge compilation, speedup learning, nonuniform complexity, memoization, amortized inference, and certifying algorithms to make that redistribution explicit. Boolean satisfiability provides a concrete setting in which construction cost, representation size, finite precision, online evaluation, and correctness can be accounted for separately. Elementary finite-domain constructions illustrate exact tabulation and neural memorization; they are background facts rather than new complexity-class collapses. Polynomial-size exact models yield a nonuniform conclusion, while uniform polynomial-time construction would yield a polynomial-time SAT algorithm. Exact caching illustrates reuse over repeated queries, and a terminating certified fallback architecture preserves correctness independently of predictor accuracy. Elementary cost bounds state sufficient conditions for expected savings and limit the overhead of failed predictions; whether learned systems meet those conditions remains an empirical and algorithmic question. The contribution is a focused resource-accounting perspective on fast and deliberative computation, with explicit assumptions for assessing when acquired knowledge repays its cost.
\end{abstract}

\section{Introduction}

A difficult decision problem may be expensive to solve on first encounter and inexpensive to answer subsequently. If an exact answer has already been computed and retained, a later query can be answered by retrieval rather than renewed search. A learned representation may extend this reuse beyond identical inputs by capturing regularities shared by different instances. These observations motivate a distinction between the cost of acquiring knowledge and the cost of using it.

Knowledge compilation provides a foundational framework: the size of a compiled representation and the tractability of the queries it supports are distinct design dimensions \citep{darwiche2002}. In the present setting, the compiled object represents the answers to a family of bounded-size SAT instances. This is a particular compilation construction, not an assertion that every conventional knowledge-compilation method stores the answers to all possible instances.

The theory of compilability formalizes how preprocessing a fixed part of an instance can simplify subsequent queries, with representation-size constraints central to the resulting complexity questions \citep{cadoli2002}. Here only the maximum input length is fixed, and the table covers all formulas up to that length. Allowing an exponentially large table gives an elementary baseline; it does not establish tractability under the succinctness constraints studied in compilability. The Learning to Reason framework likewise treats knowledge acquisition and subsequent query answering together, making the agent's environment and learning interface explicit \citep{khardon1997}.

This article is a perspective and explanatory synthesis of established ideas. The elementary proofs make resource assumptions inspectable; no new complexity separation, neural capacity bound, or empirical speedup is claimed. Building on speedup learning, knowledge compilation, and neural advice, we organize three settings for exact SAT: finite-domain restriction, access to a fixed oracle language, and model families supplied as nonuniform advice. Certification then separates correctness from savings. The contribution is this focused exposition and its explicit proof obligations, rather than priority for connecting learning to complexity classes. The organizing question is when expenditure on reusable knowledge is repaid by lower subsequent computation.

The finite-domain observation used as a starting point is precise but limited. For every fixed maximum input length, all membership answers can be encoded in a finite object. This establishes efficient online evaluation under explicit representation and access assumptions. It does not establish that the object has polynomial size, that it can be constructed in polynomial time, or that one uniform algorithm solves all input lengths efficiently. The purpose of this note is to make these distinctions explicit while retaining the useful interpretation of learning as amortization of reasoning.

The terms \emph{System 1} and \emph{System 2} are used here as a computational analogy. The former denotes rapid retrieval or evaluation of acquired knowledge, and the latter denotes additional computation used to obtain an answer not already supported by that knowledge. Neither term is identified with a complexity class.

\subsection{Closest precedents and scope}

\citet{lin2021limitations} already relate compact binary parameter strings to polynomial advice while allowing unrestricted acquisition, including oracle assistance. Their bounded-length lookup models explicitly trade exponential storage for efficient queries. This is the closest precedent for the combination of tabulation and nonuniform parameters used here. Their principal results concern weighted languages, scoring, and sampling; our decision-language exposition does not transfer those results to arbitrary SAT predictors.

The neural-advice connection is older still. \citet{balcazar1997neural} relate information in real-valued recurrent-network weights to nonuniform complexity. Their ideal real constants differ from the finite bit strings charged here. \citet{hertrich2023knapsack} explicitly distinguish uniform neural construction from nonuniform existence when discussing NP-hard knapsack. Thus neither neural advice nor the uniformity qualification is new to this perspective.

Learning to reduce future computation is the subject of speedup learning. \citet{minton1988utility} studies the utility problem: learned control knowledge can cost more to apply than it saves. \citet{tadepalli1996speedup} gives constructive learning guarantees under specified problem-solving assumptions. Offline reasoning for later efficiency is also studied by \citet{moses1996knowledge}. These precedents motivate our cost accounting; our elementary identities do not replace their learning theorems. \citet{juba2013implicit} connects learning and reasoning under PAC semantics and restricted proof systems, a different guarantee from deciding every SAT input exactly.

\citet{fortnow2022fifty} situates practical progress on hard problems alongside the unresolved worst-case question. A recent preprint by \citet{fortnow2026manage} uses $\Ppoly$-computable distributions to study machine learning; its distributional setting differs from exact SAT decision. Together these works establish the broader conceptual connection. Sections~\ref{sec:uniformity} and~\ref{sec:hybrid} specialize the accounting to exact models and certified fallback; Section~\ref{sec:open} identifies the additional evidence still needed for useful learned reasoning.

\section{Problem formulation and computational resources}
\label{sec:setup}

\subsection{The bounded SAT domain}

Fix a standard binary encoding of Boolean formulas. The unrestricted satisfiability language is
\begin{equation}
\SAT=\{\langle\varphi\rangle:\exists a\ \varphi(a)=1\}.
\end{equation}
SAT is NP-complete, and a proposed satisfying assignment can be checked in polynomial time in the encoded formula length \citep{cook1971}. Malformed encodings are assigned the answer zero. We use an encoding in which a formula of length $n$ has at most $n$ distinct variables and can be evaluated on an assignment in time polynomial in $n$.

For $N\in\mathbb{N}$, let
\begin{equation}
D_N=\bits^{\leq N},\qquad A_N=\SAT\cap D_N.
\label{eq:domain}
\end{equation}
The bound $N$ applies to the entire binary description, not merely to the number of variables or to an unencoded numerical parameter. The cardinality of the domain is
\begin{equation}
|D_N|=\sum_{n=0}^{N}2^n=2^{N+1}-1.
\end{equation}
The empty string is included and is treated as malformed. The parent problem is NP-complete; its fixed finite restriction is a different, finite language.

\subsection{Offline construction and online use}

A compiled representation and its query procedure are written as
\begin{equation}
K_N=\Compile(N),\qquad
\Answer(x;K_N)=\ind_{\SAT}(x)\quad(x\in D_N).
\label{eq:compiler}
\end{equation}
The principal resource quantities are
\begin{equation}
C_{\mathrm{offline}}(N),\qquad |K_N|,\qquad
C_{\mathrm{online}}(x;K_N).
\label{eq:resources}
\end{equation}
Correctness in \eqref{eq:compiler} means correctness on every input in the bounded domain. High test accuracy is a different guarantee and will not be substituted for it.

A small online cost does not imply a small construction cost or a small representation. In particular, a polynomial-time indexing function may address a table containing exponentially many answers. The cost of acquiring those answers remains part of the end-to-end computation, even when it is not charged to a subsequent query.

\subsection{Access model and precision}

The ordinary classes $\Pclass$ and $\NPclass$ are understood in a standard uniform Turing-machine model. A separate online-lookup analysis may assume a preinitialized, randomly addressable memory. When an array contains $2^{\Theta(N)}$ entries, its addresses require $\Theta(N)$ bits. Unit-cost word access with sufficiently long words is therefore an explicit assumption, not a consequence of a fast hash computation. Bit-level address manipulation can instead be charged polynomially in $N$.

Neural computations are evaluated using finitely represented weights and intermediate values. An exact real constant of unbounded precision is not admitted as a free repository of answers. Stochastic predictors may be used to propose candidates, but the exact-language statements below concern deterministic, total, correct decision procedures.

\section{Exact compilation on a finite domain}
\label{sec:table}

\begin{lemma}[Direct-address compilation]
\label{lem:table}
For every $N$, there is an exact representation of $A_N$ containing $2^{N+1}-1$ decision bits. Its entries can be indexed in time polynomial in the query length. With a preinitialized direct-access memory, online queries have polynomial cost; in a word-RAM with words of at least $N+1$ bits, $O(|x|+1)$ word operations suffice.
\end{lemma}

\begin{proof}
For $x\in\bits^n$, define
\begin{equation}
h(x)=2^n-1+\val(x),
\label{eq:hash}
\end{equation}
where $\val(x)$ is the binary value of $x$ and $\val(\varepsilon)=0$. For fixed $n$, the image is the interval from $2^n-1$ to $2^{n+1}-2$. These intervals are disjoint for different lengths. Thus $h$ is a bijection from $D_N$ to the array indices $0,\ldots,2^{N+1}-2$.

The binary representation of $h(x)$ can be produced using $O(n+1)$ elementary bit operations. Construct a table satisfying
\begin{equation}
K_N[h(x)]=\ind_{\SAT}(x)\qquad(x\in D_N).
\end{equation}
An exact SAT algorithm, including exhaustive assignment enumeration, computes every required entry. Online evaluation computes $h(x)$ and retrieves the corresponding bit. Under the stated direct-access convention, the retrieval is a single memory operation. Charging operations on the address at the bit level gives a polynomial bound in $N$ instead.
\end{proof}

This construction does not require the postulate that a suitable collision-free hash exists: the input supplies an efficiently computable address. The computationally difficult information is stored in the answer array. Equation~\eqref{eq:hash} is an indexing map, not a solution of SAT in disguise.

For the encoding fixed in Section~\ref{sec:setup}, exhaustive construction has the elementary upper bound
\begin{equation}
C_{\mathrm{offline}}(N)
\leq\sum_{n=0}^{N}2^n\,2^n\,\poly(n)
=O\!\left(4^N\poly(N)\right).
\label{eq:offline}
\end{equation}
This is an upper bound for a particular construction, not an optimality claim or a lower bound. Storing one satisfying assignment with every positive entry requires at most $O(N2^N)$ bits. Returning all satisfying assignments is a different task and can require exponentially long output even for a single input.

The direct-access assumption cannot be silently transferred to a sequential tape containing a long table. A sequential scan may cost time proportional to the representation length. An alternative is to hardwire the finite answer function into a decision tree or finite automaton. Such a machine reads each input symbol once, but its description may contain exponentially many states. Either implementation makes the same tradeoff visible: short response time can be purchased with a large representation or a large machine description.

\begin{figure}[!htbp]
\centering
\begin{tikzpicture}
  \node[band] at (-0.1,2.45) {Offline construction};
  \node[block,text width=3.0cm] (domain) at (1.7,1.45)
    {All inputs in $D_N$\\$|D_N|=2^{N+1}-1$};
  \node[block,text width=3.0cm] (solver) at (6.25,1.45)
    {Exact SAT computation\\$C_{\mathrm{offline}}(N)$};
  \node[memory,text width=3.0cm] (table) at (10.8,1.45)
    {Answer table $K_N$\\$K_N[h(x)]=\ind_{\SAT}(x)$};
  \draw[flow] (domain) -- (solver);
  \draw[flow] (solver) -- (table);
  \draw[black!25] (-0.1,0.37) -- (12.55,0.37);
  \node[band] at (-0.1,-0.05) {Online query};
  \node[block,text width=3.0cm] (query) at (1.7,-1.1)
    {New input $x\in D_N$};
  \node[block,text width=3.0cm] (index) at (6.25,-1.1)
    {Index $h(x)$\\Polynomial cost};
  \node[block,text width=3.0cm] (lookup) at (10.8,-1.1)
    {Direct-access lookup\\Exact answer};
  \draw[flow] (query) -- (index);
  \draw[flow] (index) -- (lookup);
  \draw[flow] (table) -- (lookup);
\end{tikzpicture}
\caption{Exact finite-domain compilation. The online computation uses previously constructed information. The direct-access memory model is explicit, and neither the offline cost nor the representation size is claimed to be polynomial in $N$.}
\label{fig:compile}
\end{figure}

\FloatBarrier
\section{Equality of finite restrictions}
\label{sec:restriction}

\begin{definition}[Finite-domain restriction of a language class]
For a class $\mathcal C$ of languages, define
\begin{equation}
\mathcal C\restr D_N=\{L\cap D_N:L\in\mathcal C\}.
\end{equation}
The notation $2^{D_N}$ denotes the power set of $D_N$.
\end{definition}

\begin{lemma}[Equality on every fixed finite domain]
\label{lem:restriction}
For every fixed $N$,
\begin{equation}
\boxed{\Pclass\restr D_N=\NPclass\restr D_N=2^{D_N}.}
\label{eq:finite-equality}
\end{equation}
\end{lemma}

\begin{proof}
Let $S\subseteq D_N$. Since $S$ is finite, it can be represented by a finite trie of its strings, with accepting states marking exactly the strings in $S$. A fixed finite automaton implements this trie, rejecting inputs outside $S$. On an input of length $n$, it takes $O(n+1)$ steps. Hence $S\in\Pclass\subseteq\NPclass$.

Every subset of $D_N$ consequently occurs as the restriction of a language in $\Pclass$, and therefore also as the restriction of a language in $\NPclass$. The reverse inclusions in $2^{D_N}$ hold by definition.
\end{proof}

The underlying standard fact is that every finite language is regular; see \citet[Remark~1.12, p.~12]{rigo2014}. Lemma~\ref{lem:restriction} is an immediate corollary, not a SAT-specific tractability result. Writing $\mathsf{REG}$ for the regular languages and using $\mathsf{REG}\subseteq\Pclass\subseteq\NPclass$, the same argument gives
\begin{equation}
\mathsf{REG}\restr D_N=\Pclass\restr D_N=\NPclass\restr D_N=2^{D_N}.
\end{equation}
The equality is not specific to SAT or to learning. It records the fact that asymptotic complexity distinctions cannot be inferred from a single finite domain. Its operational interpretation is nevertheless useful: for a prescribed domain, exact answers can be made available through an appropriate compiled representation.

The quantifiers are essential. The statement that
\begin{equation}
\forall N\ \exists K_N
\end{equation}
with efficient access for that fixed $N$ does not supply one procedure that constructs all $K_N$ and answers all input lengths within a common polynomial bound. The description size, construction cost, and implementation constants may depend arbitrarily on $N$.

\begin{corollary}
A fixed finite restriction $A_N$ is in $\Pclass$. If such a restriction were NP-hard under polynomial-time many-one reductions, then $\Pclass=\NPclass$.
\end{corollary}

\begin{proof}
The first assertion follows from Lemma~\ref{lem:restriction}. For the second, compose a polynomial-time reduction from any language in $\NPclass$ to $A_N$ with the polynomial-time decision procedure for $A_N$.
\end{proof}

Thus the NP-completeness of the unrestricted parent problem must not be attributed automatically to its fixed finite restriction.

\section{The distinction from relativized complexity}
\label{sec:oracles}

For a fixed language $A\subseteq\bits^*$, let $\Pclass^A$ and $\NPclass^A$ denote the deterministic and nondeterministic polynomial-time oracle classes. The oracle gives an exact membership bit; writing a query and processing the answer are charged, while the oracle's internal decision procedure is not. A Karp reduction is a polynomial-time many-one map satisfying $x\in L\iff f(x)\in A$. A Cook reduction is a polynomial-time oracle computation and may make multiple adaptive queries. These are different reduction notions, neither of which makes the target language polynomial-time decidable by definition.

\citet{baker1975} established the existence of fixed oracles $A$ and $B$ with
\begin{equation}
\Pclass^A=\NPclass^A,\qquad
\Pclass^B\neq\NPclass^B.
\end{equation}
These statements concern computations on inputs of unbounded length. They do not follow simply from the availability of a finite table. The relevant distinction can be shown directly.

\begin{lemma}[Elimination of a polynomial-time oracle]
\label{lem:eliminate}
If $A\in\Pclass$, then
\begin{equation}
\Pclass^A=\Pclass,\qquad\NPclass^A=\NPclass.
\label{eq:oracle-eliminate}
\end{equation}
\end{lemma}

\begin{proof}
Let an oracle machine have polynomial time bound $p(n)$. It makes at most $p(n)$ queries and each has length at most $p(n)$. Let $q(s)$ be a nondecreasing polynomial time bound for an ordinary decider for $A$. Replace every oracle query by that decider. The total simulation time is at most
\begin{equation}
O\!\left(p(n)+p(n)q(p(n))\right),
\end{equation}
which is polynomial. This proves $\Pclass^A\subseteq\Pclass$. For a nondeterministic machine, the same replacement is performed along each branch and preserves acceptance. Therefore $\NPclass^A\subseteq\NPclass$. The reverse inclusions hold because an oracle can be ignored.
\end{proof}

Applying this lemma to $A_N$ gives
\begin{equation}
\Pclass^{A_N}=\Pclass,\qquad
\NPclass^{A_N}=\NPclass,
\end{equation}
and consequently
\begin{equation}
\Pclass^{A_N}=\NPclass^{A_N}
\quad\Longleftrightarrow\quad
\Pclass=\NPclass.
\end{equation}
There is no contradiction with \eqref{eq:finite-equality}. Restricting the inputs of an entire problem to $D_N$ is different from giving an otherwise unrestricted machine access to a fixed finite oracle.

The same argument applies to a uniformly polynomial-time learned predictor through its actual output language. If $h$ is a total deterministic polynomial-time algorithm and $A_h=\{x:h(x)=1\}$, then $\Pclass^{A_h}=\Pclass$ and $\NPclass^{A_h}=\NPclass$. This does not imply that $A_h$ equals the intended target language. Target correctness is an additional property.

\section{Neural representations and their costs}
\label{sec:neural}

\subsection{Inference and prescribed training}

Let a feedforward network have widths $d_0,\ldots,d_\ell$, and put
\begin{equation}
S=\sum_{i=1}^{\ell}(d_{i-1}d_i+d_i).
\end{equation}
Assume conventional dense arithmetic and activations whose evaluation and differentiation have bounded elementary cost. A forward pass has arithmetic cost $O(S)$. Backpropagation performs matrix-vector or matrix-matrix operations of the same asymptotic size. For $m$ already-labelled examples and $t$ prescribed epochs, ordinary gradient-based training therefore has the operation counts
\begin{equation}
T_{\mathrm{infer}}=O(S),\qquad
T_{\mathrm{train}}=O(tmS).
\label{eq:arithmetic}
\end{equation}
Parameter updates of linear cost in the number of parameters are included in these bounds.

If layer widths are $O(N)$, then
\begin{equation}
T_{\mathrm{infer}}=O(\ell N^2),\qquad
T_{\mathrm{train}}=O(tm\ell N^2).
\label{eq:dense}
\end{equation}
These expressions are not architecture-independent formulas. With $b$-bit intermediate values and numerical-operation cost at most $a(b)=\poly(b)$, the corresponding bit-cost bounds are $O(Sa(b))$ and $O(tmSa(b))$. They are polynomial in the original instance length only if network size, precision, sample count, and iteration count are themselves appropriately bounded.

A prescribed training budget bounds computation performed, not the quality achieved. The existence of an exact model does not prove that a given optimization method will find it within that budget. For example, \citet{blum1988} proved NP-completeness of deciding whether a particular small threshold network can fit supplied examples exactly. This is a fitting decision problem, not a claim that one gradient update has superpolynomial cost. Obtaining exact training labels must also be charged separately when those labels are not already available.

\subsection{Recurrent execution}

For a recurrent model executed for $r$ steps, with per-step cost $S_{\mathrm{step}}$, the analogous arithmetic bounds are
\begin{equation}
T_{\mathrm{infer}}=O(rS_{\mathrm{step}}),\qquad
T_{\mathrm{train}}=O(tmrS_{\mathrm{step}}),
\end{equation}
under the corresponding assumptions for backpropagation through time. Polynomial execution requires a polynomial bound on $r$, as well as on precision and the per-step computation. A recurrent architecture by itself supplies no such bound. Rational-weight recurrent networks can simulate Turing machines \citep{siegelmann1995}. Accordingly, the relevant claim is that \emph{bounded, finite-precision neural computations with polynomial resources are polynomial-time simulable}, not that every neural-network algorithm has polynomial total runtime.

Generated reasoning steps must likewise be charged. \citet{merrill2024cot} relate chain-of-thought length to transformer expressivity under explicit architectural assumptions. For a deployed decoder, each additional token incurs computation, and its cost may increase with the growing context. One must sum these per-step costs; a polynomial bound for a single forward pass does not bound an unbounded generation process.

\subsection{Exactness on the restricted domain}

Replacing $K_N$ by a model $f_{\theta_N}$ gives the postulate
\begin{equation}
f_{\theta_N}(x)=\ind_{\SAT}(x)
\qquad\text{for all }x\in D_N.
\label{eq:exact-model}
\end{equation}
At one fixed $N$, such a representation exists without a learning theorem. The meaningful further question is whether it is compact, efficiently evaluable, and obtainable at acceptable cost. A successful empirical predictor can be useful without satisfying \eqref{eq:exact-model}, but then it is not an exact SAT oracle.

\section{An explicit exact ReLU construction}
\label{sec:relu}

\begin{lemma}[Neural realization of a finite answer table]
\label{lem:relu}
For every fixed $N$, the SAT decision function on exactly $N$ input bits has an exact one-hidden-layer ReLU representation with at most $2^N$ hidden units and $O((N+1)2^N)$ parameters. With a fixed dense serialization convention, architecture and parameters admit an $O((N+1)2^N)$-bit description. Evaluation on binary inputs can use exact integer arithmetic with $O(\log(N+2))$ bits per numerical intermediate.
\end{lemma}

\begin{proof}
For each $u\in\bits^N$, define
\begin{equation}
r_u(x)=\ReLU\!\left(
1-\sum_{j:u_j=0}x_j-\sum_{j:u_j=1}(1-x_j)
\right).
\label{eq:relu-unit}
\end{equation}
For a binary input $x$, the two sums together equal the Hamming distance between $x$ and $u$. Therefore
\begin{equation}
r_u(x)=\begin{cases}1,&x=u,\\0,&x\neq u.\end{cases}
\end{equation}
Define the linear readout
\begin{equation}
f_N(x)=\sum_{u\in\SAT\cap\bits^N}r_u(x).
\label{eq:relu-sat}
\end{equation}
Exactly one indicator $r_u$ is nonzero at each binary input if all indicators are present; after selecting the accepting indices, the sum is one precisely on SAT encodings. Each unit has $N$ input weights and a bias. The unit weights are $+1$ or $-1$, and the bias is $1-|u|_1$, of magnitude at most $\max\{1,N\}$. The parameter-count bound follows. An empty sum is implemented as the constant zero function.

For the description bound, serialize the dense layers in a fixed order, encoding each input weight by one bit, each bias by $O(\log(N+2))$ bits, and each readout weight (equal to one) by constant space. The dimensions require only $O(N+\log(N+2))$ header bits. For at most $2^N$ hidden units this totals $O((N+\log(N+2))2^N)=O((N+1)2^N)$ bits. Each hidden preactivation and each sequential partial sum used to compute it has magnitude $O(N+1)$. On binary inputs at most one hidden output is one, so every partial readout sum is zero or one. Thus $O(\log(N+2))$ bits suffice for each numerical intermediate; loop indices and addresses are accounted for separately.
\end{proof}

Finite-sample memorization is a well-established expressivity question; for example, \citet{yun2019} analyze capacity bounds for ReLU networks with different depth and width constraints. The construction here is the elementary Boolean-cube indicator construction and makes no optimality claim.

This construction is explicit and exact on binary inputs, but it is not a polynomial-size construction in $N$. Evaluating all hidden units by a conventional dense forward pass can require $O(N2^N)$ arithmetic operations. The accepting units themselves are selected using the already-known SAT answers. Consequently, the construction does not resolve the cost of acquiring those answers.

For lengths up to $N$, one may select a length-specific subnetwork or use an injective encoding that includes the length. The conclusion remains the same: finite representability is automatic, whereas useful compression is a substantive additional property.

A small-range hash does not evade this issue. If a retrieval rule depends only on $h(x)$, correctness requires all inputs mapped to the same code to have the same answer. An injective map guarantees this by using many addresses. A label-preserving compressed map may use fewer addresses, but constructing or evaluating that map can contain the original decision difficulty. In the extreme case $h(x)=\ind_{\SAT}(x)$, the hash is already the SAT decision algorithm.

\section{Scaling, advice, and uniform construction}
\label{sec:uniformity}

The fixed-size observation becomes a complexity-theoretic hypothesis only when a family of sizes is considered. Let $\theta_n$ encode a model for inputs in $\bits^n$. Its description includes both architecture and finitely represented parameters.

\begin{proposition}[Compact models and uniform models]
\label{prop:advice}
Suppose there is a polynomial $q$ and one deterministic evaluation algorithm $E$ such that, for every $n$, a parameter string $\theta_n$ satisfies $|\theta_n|\leq q(n)$ and
\begin{equation}
E(x,\theta_n)=\ind_{\SAT}(x)\qquad(x\in\bits^n),
\end{equation}
with evaluation time polynomial in $n$. Then $\SAT\in\Ppoly$. If, in addition, one deterministic algorithm constructs $\theta_n$ from $1^n$ in polynomial time, then $\SAT\in\Pclass$ and $\Pclass=\NPclass$.
\end{proposition}

\begin{proof}
The first assertion is the definition of polynomial-time computation with polynomial-length advice: $\theta_n$ is the advice string for input length $n$, and $E$ is the fixed evaluator. No uniform construction of the advice is required.

For the second assertion, given $x$ of length $n$, construct $\theta_n$ and then run $E(x,\theta_n)$. Both stages take polynomial time, so this is a uniform polynomial-time decider for SAT. Composing it with the polynomial-time reductions supplied by NP-completeness puts every language in $\NPclass$ in $\Pclass$.
\end{proof}

The first conclusion is an application of the standard advice framework \citep{karp1980}. It also has a stronger known consequence: $\SAT\in\Ppoly$ implies $\NPclass\subseteq\Ppoly$. To see this, a reduction from length-$n$ inputs produces SAT instances of lengths at most $p(n)$ for a polynomial $p$; concatenating the advice strings for all lengths up to $p(n)$ still uses polynomial space. The Karp--Lipton theorem, in its second-level form incorporating Sipser's improvement, then yields
\begin{equation}
\mathsf{PH}=\Sigma_2^{\Pclass},
\end{equation}
where $\mathsf{PH}$ is the polynomial hierarchy \citep{karp1982}. Thus compact exact model families are a substantial complexity-theoretic hypothesis even without an efficient construction algorithm. This consequence does not by itself assert $\Pclass=\NPclass$.

The hypothesis includes correctness on all inputs, not merely on a training distribution. If construction is described as learning, the uniform polynomial budget must include data acquisition, exact label generation when required, parameter fitting, and any other computation needed to obtain the exact model. Unlimited preprocessing cannot be converted into a uniform polynomial-time algorithm by charging only inference.

The model for length $n$ must be shared by \emph{all} length-$n$ inputs. If a different parameter string is supplied for each individual input, that string could simply contain its answer; this is not length-dependent advice. A randomized training procedure also does not automatically satisfy the deterministic construction hypothesis. Its success probability, certification procedure, and cost of failed runs require a separate analysis. In particular, a small measured test error does not prove correctness simultaneously on all $2^n$ inputs.

The distinction has direct precedents in \citet{lin2021limitations} and \citet{hertrich2023knapsack}. In particular, assuming only $\Pclass\neq\NPclass$ does not exclude polynomial-size nonuniform SAT models. The stronger assumption $\NPclass\nsubseteq\Ppoly$ would exclude the exact, polynomially evaluable families in Proposition~\ref{prop:advice}. Neither assumption is proved here.

\paragraph{Decision versus witness production.}
An exact SAT decider can recover one satisfying assignment by self-reduction: after a positive decision, tentatively fix each variable to zero, query the restricted formula, and choose zero if satisfiability is retained and one otherwise. Each choice preserves a satisfying extension, so after at most $n$ additional queries an assignment is obtained. With a standard encoding the restricted formulas have length at most a polynomial $p(n)$, and constructing them costs polynomial time. Thus uniform polynomial-time decision also gives polynomial-time witness production. A finite table only supports this reduction if it covers \emph{all} intermediate encodings; coverage of the original input alone is insufficient. For nonuniform families, concatenating the advice for lengths up to $p(n)$ still uses polynomial space. This supplies one witness, not an enumeration of all witnesses.

A fixed model for each length and a model retrained after every query also require different descriptions. The latter is a stateful procedure whose answers may depend on query history. It is not automatically a fixed oracle language. Nevertheless, if its state has polynomial size and it performs polynomially many polynomial-cost updates, those updates can be simulated within polynomial time. Any external exact-solving costs must be added rather than hidden inside the update operation.

\section{Amortization by exact caching}
\label{sec:cache}

Memoization stores the results of earlier computations for later reuse, a connection to machine learning articulated by \citet{michie1968}. The simplest System 1/System 2 implementation is an exact dictionary. On a previously solved input, the fast component retrieves the answer. On a cache miss, a complete solver computes the answer and the result is stored. The solver is an ordinary deterministic algorithm; its designation as System 2 supplies no nondeterministic computational power.

\begin{lemma}[Amortized cost of exact reuse]
\label{lem:cache}
Let $x_1,\ldots,x_Q\in D_N$ be a query sequence. Assume an initially empty cache retains every exact answer without eviction. Let $c_N$ bound the per-query lookup, bookkeeping, and insertion cost, and let $s(x)$ be the runtime of the complete solver on a previously uncached input $x$. If
\begin{equation}
U_Q=\{x_i:1\leq i\leq Q\},
\end{equation}
then the total runtime satisfies
\begin{equation}
\boxed{T_Q\leq Qc_N+\sum_{x\in U_Q}s(x).}
\label{eq:cache-bound}
\end{equation}
\end{lemma}

\begin{proof}
Every query incurs at most $c_N$ dictionary overhead. Each distinct input invokes the solver only at its first occurrence. The answer is then retained, so all later occurrences use retrieval. Summing the two types of cost gives \eqref{eq:cache-bound}.
\end{proof}

Because the domain is finite and the solver terminates on every input,
\begin{equation}
s_N=\max_{x\in D_N}s(x)<\infty.
\end{equation}
It follows that
\begin{equation}
\frac{T_Q}{Q}\leq c_N+\frac{|U_Q|s_N}{Q}
\leq c_N+\frac{|D_N|s_N}{Q}.
\label{eq:amortized}
\end{equation}
For fixed $N$, the upper bound on the amortized solving contribution tends to zero as $Q$ grows. This statement is a limit in the number of queries at fixed domain size, not a polynomial bound as $N$ grows.

Under a random-access implementation, a trie or dictionary can provide polynomial per-query overhead in $N$ while accumulating up to exponentially many entries. The required storage must therefore remain explicit. Before repeats occur, an adversarial query stream can repeatedly invoke the solver. Cache eviction, memory limits, or changes in the task can reintroduce earlier costs; the no-eviction assumption is part of the lemma, not a guarantee about practical learning systems.

For independent queries with common probabilities $p_x=\Prb(X=x)$, the same argument gives a finite-workload bound:
\begin{equation}
\E[T_Q]\leq Qc_N+
\sum_{x\in D_N}s(x)\bigl[1-(1-p_x)^Q\bigr].
\label{eq:cache-distribution}
\end{equation}
Indeed, the bracket is the probability that $x$ occurs at least once, and each such occurrence triggers exactly one solve. Linearity of expectation then applies; the occurrence indicators for different strings need not be independent. For a uniform distribution over $M$ inputs, the expected number of solves is $M[1-(1-1/M)^Q]$. When $Q$ is small compared with $M$, this is close to $Q$. The fixed-domain limit therefore need not predict useful savings at the query volumes actually encountered.

\section{Certified learning and adaptive fallback}
\label{sec:hybrid}

\subsection{From memorization to structural reuse}

A learned predictor can reuse information across distinct inputs rather than only across repeated encodings. For SAT, a concrete candidate model represents a formula by a graph linking clauses and literals, applies message passing, and produces an assignment or a satisfiability prediction. NeuroSAT provides a research example of a message-passing model trained from satisfiability labels that can recover satisfying assignments on studied instances \citep{selsam2019}. Its empirical generalization is not an all-instance polynomial-time SAT guarantee.

For a proposed assignment $\widehat a$, the predicate
\begin{equation}
V(\varphi,\widehat a)=\ind\{\varphi(\widehat a)=1\}
\end{equation}
certifies a positive answer when it evaluates to one. A failed check does not certify unsatisfiability. A model's confident negative prediction is likewise not a proof merely because its confidence is high.

An exact hybrid can return a previously cached exact answer, accept a satisfying assignment after verification, or invoke a complete solver. A negative result may also be accepted using a sound checked refutation; generating and checking the refutation are part of the cost, and no polynomial upper bound on its length is assumed. The fallback criterion must be implementable. There is no free mechanism that recognizes all erroneous model predictions.

This separation between an answer-producing procedure and a trusted checker follows the established principle of certifying algorithms \citep{mcconnell2011}. Neural guidance can also be integrated within exact search: NeuroCore uses learned unsatisfiable-core predictions to guide branching in existing SAT solvers \citep{selsam2019core}. Such guidance differs from the separate fast path considered here, but likewise separates prediction quality from logical correctness.

An earlier practical precedent is SATzilla, which learns to select a solver for an instance \citep{xu2008satzilla}. Selection overhead must be repaid by reduced solving time. Selection and internal search guidance can change the fallback runtime itself; the comparison below identifies explicitly when it assumes an unchanged baseline solver.

For a refutation $\pi$, a checker polynomial in $|\varphi|+|\pi|$ need not run in time polynomial in $|\varphi|$ alone. In the Cook--Reckhow framework, the existence of a sound, complete, polynomially checkable propositional proof system with polynomial-size proofs for every tautology is equivalent to $\NPclass=\coNP$ \citep{cook1979}. The corresponding statement for UNSAT follows by negation. The architecture assumes neither universally short negative certificates nor an efficient procedure for finding them.

\begin{proposition}[Soundness of certified fallback]
Assume the fast-path procedure terminates on every input, returning either a certified answer or a fallback signal. If each returned answer is justified by an exact cache entry or a sound verifier for that answer, and each fallback signal invokes a terminating exact solver, then the hybrid is a total correct SAT decider on its domain, irrespective of the predictor's error rate.
\end{proposition}

\begin{proof}
A fast-path answer is correct by its certification condition. Any input not answered on that path reaches the exact solver and is answered correctly there. Fast-path termination ensures that one of these cases is reached on every input. Solver termination then guarantees termination of the entire decision procedure.
\end{proof}

A finite timeout covering prediction, certificate generation, and checking can enforce the fast-path termination condition. An unfinished or failed check must trigger fallback. Soundness alone is insufficient: a predictor that loops forever can otherwise prevent the complete solver from being reached.

\begin{figure}[tbp]
\centering
\begin{tikzpicture}
  \node[block,text width=2.55cm] (query) at (1.5,2.15)
    {SAT input $x$};
  \node[memory,text width=3.4cm] (fast) at (6.0,2.15)
    {System 1\\Exact cache or learned model};
  \node[block,text width=2.75cm] (gate) at (10.75,2.15)
    {Sound check\\or exact cache};
  \node[block,text width=2.75cm] (answer) at (10.75,-0.1)
    {Exact answer};
  \node[block,text width=3.4cm] (slow) at (6.0,-0.1)
    {System 2\\Complete SAT solver};
  \node[memory,text width=3.4cm] (learn) at (6.0,-2.4)
    {Store exact result\\Optional learned update};
  \draw[flow] (query) -- (fast);
  \draw[flow] (fast) -- (gate);
  \draw[flow] (gate) -- node[note,right] {certified} (answer);
  \draw[flow] (gate.south west) -- node[note,above,sloped] {otherwise} (slow.north east);
  \draw[flow] (slow) -- (answer);
  \draw[flow] (slow) -- (learn);
  \draw[optional] (learn.west) -- (3.55,-2.4) -- (3.55,0.85) -- (fast.south west);
  \node[note,text width=2.55cm] at (1.7,-1.15)
    {Reuse reduces\\later search};
\end{tikzpicture}
\caption{An exact learned-prediction-and-fallback architecture. The acceptance test is a real computation, not an omniscient error detector. Verified positive witnesses and exact cached results can bypass further search. Uncertified queries use the complete solver. The dashed route makes learned updating optional; it does not make exact solving free.}
\label{fig:hybrid}
\end{figure}

\subsection{Total and expected costs}

For a realized query sequence, let $F\subseteq\{1,\ldots,Q\}$ index the queries that invoke the complete solver. Group every operation into initial construction, fast-path processing, exact solving, or updating. With this accounting,
\begin{equation}
T_Q=C_{\mathrm{initial}}+
\sum_{i=1}^{Q}C_{\mathrm{fast},i}+
\sum_{i\in F}\bigl(C_{\mathrm{solve},i}+C_{\mathrm{update},i}\bigr).
\label{eq:hybrid-cost}
\end{equation}
The fast-path term includes model evaluation, verification, and the decision to fall back. The update term includes optional fitting and associated bookkeeping. A chosen implementation may also perform updates on certified fast-path outcomes; those costs must then be added explicitly or included in the corresponding fast-path term.

Once an exact answer has been obtained, retraining can be omitted. This sets the optional fitting contribution to zero but leaves the solving cost unchanged. Retraining may still be worthwhile if it reduces later fallback frequency. Polynomially many polynomial-time updates add only polynomial work, but they do not establish a polynomial bound for the computations used to obtain exact labels.

For a fixed model queried on a distribution $\mathcal D_N$, let $\rho_N$ be the probability of requiring fallback. A one-query expectation, conditional on the current stored state, is
\begin{equation}
\begin{split}
\E[T(X)]={}&\E[C_{\mathrm{fast}}(X)]\\
&+\rho_N\E\!\left[
C_{\mathrm{solve}}(X)+C_{\mathrm{update}}(X)
\mid\text{fallback}\right].
\end{split}
\label{eq:expected}
\end{equation}
When $\rho_N=0$, the second term is interpreted as zero. For an adaptive model, the same identity can be applied conditionally on the query history and then averaged over that history. A low fallback rate can yield a substantial expected speedup without any polynomial worst-case guarantee. Rare expensive cases must still be counted.

These identities account for work; they do not prove that the hybrid is faster. For a specified solver-only baseline with cost $B_i$ on query $i$, end-to-end improvement on the same sequence requires
\begin{equation}
C_{\mathrm{initial}}+\sum_{i=1}^{Q}C_{\mathrm{fast},i}
+\sum_{i\in F}(C_{\mathrm{solve},i}+C_{\mathrm{update},i})
<\sum_{i=1}^{Q}B_i.
\label{eq:break-even}
\end{equation}
For comparison with a cache-only baseline, its actual lookup, storage, and miss-solving costs must replace the right-hand side. All terms must use the same cost model. Verification overhead, training expenditure, or a concentration of hard instances among fallbacks can prevent improvement even when prediction accuracy is high. The following elementary consequences make sufficient assumptions explicit; they do not show that a trained system meets them.

\subsection{A sufficient condition for savings and a bound on overhead}

\begin{proposition}[Expected break-even for a fixed policy]
\label{prop:savings}
Fix a query distribution on $D_N$, a deterministic baseline solver with runtime $s(x)$, and a deterministic certified fast-path policy that remains unchanged across queries. On fallback, run that same solver from its usual initial state. Let $H(x)$ indicate certified fast-path success, and let $f(x)$ include all fast-path and controller costs. There are no online updates. Put
\begin{equation}
\mu=\E[s(X)],\qquad
\delta=\E[H(X)s(X)]-\E[f(X)].
\end{equation}
For initial acquisition cost $C$ and $Q$ queries with the specified marginal distribution,
\begin{equation}
\E[T_Q]=C+Q(\mu-\delta).
\label{eq:savings}
\end{equation}
If $\delta>0$, the hybrid beats the solver-only expected total $Q\mu$ whenever $Q>C/\delta$.
\end{proposition}

\begin{proof}
The online cost on $x$ is $f(x)+(1-H(x))s(x)$. Taking its expectation gives $\mu-\delta$. Summing over $Q$ queries and adding $C$ proves the identity and the threshold. Independence between queries is unnecessary because the policy is fixed and only the marginal expectations are used.
\end{proof}

The saved quantity is solver time on certified successes, not raw accuracy. For example, if $\Prb(H=1)\geq\alpha\geq0$, every success has $s(x)\geq b\geq0$, and $\E[f(X)]\leq a$, then $\delta\geq\alpha b-a$. Thus $\alpha b>a$ is a sufficient condition. Conversely, with probability $0.99$ let the solver cost one unit and the fast path succeed; with probability $0.01$ let the solver cost $10^6$ units and the fast path fail. If the fast path always costs two units, then $\delta=0.99-2<0$: even a 99\% certified success rate increases expected online cost. This is an illustrative workload, not an empirical measurement.

\begin{proposition}[Bounded overhead with a budgeted fast path]
\label{prop:overhead}
Suppose every fast path, including verification, timeout handling, and dispatch, costs at most $\tau(n)$ on an input of length $n$. Any updates cost at most $u(n)$, and fallback runs the unchanged baseline solver in cost $s(x)$. Then the total charged work obeys
\begin{equation}
T_Q\leq C_{\mathrm{initial}}+
\sum_{i=1}^Q\bigl[s(x_i)+\tau(|x_i|)+u(|x_i|)\bigr].
\label{eq:overhead}
\end{equation}
\end{proposition}

\begin{proof}
On a certified success the solver cost is zero; otherwise it is $s(x_i)$. Both cases cost at most $s(x_i)+\tau(|x_i|)+u(|x_i|)$. Summing proves the bound.
\end{proof}

This is an additive overhead guarantee, not a speedup or a universal multiplicative guarantee. It requires that the fast path and updates leave the baseline computation and its resources available; any initialization, interference, or copying cost must be included. It concerns charged work in the specified model, not an automatic wall-clock guarantee on shared hardware. Polynomial budgets for the fast path and updates still leave the potentially expensive term $s(x)$ intact.

\section{System 1, System 2, and the reuse of computation}
\label{sec:discussion}

The finite table represents an extreme form of acquired knowledge. After all relevant answers have been computed, the online phase needs only recognition and retrieval. Exact caching acquires that knowledge incrementally, while learning attempts a more economical representation that can transfer information across inputs. The common mechanism is reuse of earlier computational effort.

Amortized inference also has algorithmic precedents: \citet{stuhlmuller2013} invest computation in learning stochastic inverses to support subsequent inference in Bayesian networks. Their probabilistic setting differs from exact SAT decision, but directly illustrates the acquisition-versus-use distinction.

This interpretation has a cognitive counterpart. \citet{gershman2014} study amortized inference, in which earlier sub-inferences are reused across related queries, and report experimental evidence consistent with such reuse. It is an account of interactions between inference and memory, not a demonstration that human cognition solves unrestricted NP-complete problems in polynomial time.

Reliable fast judgment also depends on the task environment. \citet{kahneman2009} emphasize environmental predictability and opportunities to learn its regularities when assessing intuitive expertise. In the present framework, these considerations motivate examining the query distribution and the scope over which stored or learned knowledge remains applicable. Subjective confidence is not used as an exactness guarantee.

A direct machine-learning analogue is investigated by \citet{yu2024}: outputs produced using higher-compute reasoning procedures are distilled into models that answer without reproducing the same intermediate reasoning sequence. Their reported successes for several techniques provide empirical support for transferring some online computation into parameters. They do not establish exactness on all inputs or eliminate the need for additional reasoning on every task. We cite this work for the distillation mechanism, not for an oracle-machine analysis or a connection to $\Pclass$ versus $\NPclass$. The connection pursued in this perspective is to ask what such a distilled representation would imply under progressively stronger assumptions: exactness on a fixed finite domain, polynomial-time evaluation with polynomial-size advice across lengths, and uniform polynomial-time acquisition. These assumptions lead to different conclusions, as Sections~\ref{sec:restriction}, \ref{sec:oracles}, and~\ref{sec:uniformity} explain.

The appropriate resource description is therefore joint rather than one-dimensional:
\begin{equation}
\mathcal R_N=
\left(
|K_N|,\ C_{\mathrm{acquire}}(N),\ C_{\mathrm{online}},\
\rho_N,\ \varepsilon_N
\right),
\end{equation}
where $\rho_N$ is fallback frequency and $\varepsilon_N$ is the error rate under a specified evaluation criterion. For the certified exact hybrid, $\varepsilon_N=0$ under the stated assumptions even when the predictor alone is inaccurate. Its fallback cost may nevertheless be large.

An exhaustive table offers exactness at high storage and acquisition cost. A compact predictor may offer economical reuse but needs an explicit correctness or error model. A solver-predictor combination may lower average response cost when saved search exceeds its acquisition and operating overhead. This motivates studying when the expense of acquiring reusable knowledge is repaid by reduced future search.

Topological compactness and sparse training coverage alone do not determine any of these complexity conclusions. Here the bounded domain is finite because its binary encoding length is bounded. A previously unseen input need not be difficult, a sparse training set need not force retraining on every query, and failure of one predictor is not a lower bound against every algorithm. These distinctions keep the System 1/System 2 interpretation separate from unsupported complexity-class claims.

\subsection{Lessons from claimed complexity proofs}

\citet{deolalikar2010claim} circulated a proposed proof of $\Pclass\neq\NPclass$ connecting least fixed-point logic with the structure of SAT solution spaces. It is a failed separation claim, not a learning-based proof of $\Pclass=\NPclass$. In a technical letter published by Richard Lipton, \citet{immerman2010objections} identifies two defects: an order-independent final answer does not force intermediate fixed-point computations to be order-independent, and restricting to monadic fixed points omits relations needed to preserve the claimed computational power. The resulting restricted formalism does not capture all of $\Pclass$. The relevant lesson is to justify the passage from a chosen representation to claims about every algorithm. Here a large table or network is only an upper-bound construction; its size is not a lower bound on other SAT representations. Conversely, a SAT formula already describes its satisfying assignments compactly, without providing an efficient procedure for deciding whether any exist.

\citet{wen2010mirrored} explicitly links learned recognition, brain-inspired representations, oracle machines, and a claimed $\Pclass=\NPclass$ conclusion. The manuscript does not provide a uniform polynomial-time SAT decider or the necessary acquisition and evaluation bounds. Its oracle terminology also fails to distinguish oracle access from nondeterministic computation. The subsequent argument of \citet{wen2010recognition} identifies deterministic and nondeterministic recognition of a language with equality of the classes, without establishing polynomial runtime and NP-hardness. Even showing one language belongs to $\Pclass\cap\NPclass$ would not imply $\Pclass=\NPclass$. These are our diagnoses of the arguments, not an attribution of a published external refutation. Their thematic proximity makes the missing obligations instructive: specify the recognized language, the reduction from SAT if needed, the uniform construction, and the full runtime bound. Proposition~\ref{prop:advice} states sufficient assumptions; naming learned recognition an ``oracle'' does not establish them.

\section{Remaining questions and an evaluation protocol}
\label{sec:open}

The preceding results close logical gaps in the exposition but do not supply a learning algorithm or performance measurements. The following are unresolved obligations for the systems considered here, not claims that the wider literature has never studied them.

\paragraph{Compactness and acquisition.}
For a specified family of SAT instances, can a representation smaller than a table be constructed within a stated budget and evaluated with a useful correctness guarantee? The input family, coverage, precision, and acquisition interface must be fixed. Tractable subclasses and distributional guarantees can permit positive answers without meeting the all-input hypotheses of Proposition~\ref{prop:advice}. For unrestricted SAT, this paper establishes neither polynomial-size exact models nor lower bounds excluding them. In particular, the exponential size of our ReLU construction proves neither necessity nor efficient learnability of a smaller alternative.

\paragraph{Generalization, memory, and changing workloads.}
When does a learned representation save more work than storing previous answers under the same memory budget? Equation~\eqref{eq:cache-distribution} supplies a no-eviction reference, but bounded-memory caching, eviction, retraining, and distribution shifts change the accounting. The fixed-policy threshold in Proposition~\ref{prop:savings} does not apply unchanged to adaptation; conditional savings must be evaluated at each history, including update costs. Certification preserves correctness under a changed workload within the covered domain, but does not preserve a favorable fallback rate or expected cost.

\paragraph{Certification and deployment budgets.}
Which inputs admit certificates that the fast path can actually produce and check within its budget? This is stronger than the existence of a short certificate. Positive and negative outcomes should be reported separately because their verification requirements differ. A timeout ensures that the fallback is reached; Proposition~\ref{prop:overhead} bounds additional work only when the solver and update-budget assumptions are implemented. Resource limits that also abort the fallback yield an incomplete service returning ``unknown'', not the total decider proved above.

\paragraph{A reproducible comparison.}
An empirical study should fix encodings, instance generators, query order, memory limits, solver versions, and hardware before comparing a solver-only baseline, an exact cache, and a certified learned hybrid. Include repeat-heavy streams, distinct held-out formulas, and shifted workloads, with satisfiable and unsatisfiable instances represented. Separate identical-input reuse from transfer to distinct inputs; freeze a test stream and record any online learning chronologically to avoid using later answers earlier. For each method report acquisition and label-generation costs, parameter and cache bytes, verification work, certified success rates, conditional fallback costs, updates, and total cost as a function of query count. Report both latency distributions and aggregate work; concurrent hardware use must not disappear from the latter. Timeouts and unsolved cases must remain in the accounting. If acquisition cost is unavailable, an online-only comparison cannot establish end-to-end break-even. No measurements are supplied here; this protocol specifies how the perspective could be tested.

\section{Conclusion}

A resource-accounting perspective on learned reasoning separates the expense of acquiring knowledge from the expense of storing, evaluating, and certifying it. Connecting this distinction to oracle machines and the P versus NP question requires specifying whether the representation is fixed, supplied separately for each input length, or constructed uniformly. The established theories of compilation, nonuniform computation, memoization, and amortized inference provide complementary ways to describe this transfer of work across time and across related queries.

Finite-domain tabulation and neural memorization supply transparent baseline constructions. Their existence is elementary and does not establish compactness, efficient acquisition, or unrestricted tractability. Polynomial-size exact model families have the standard nonuniform consequences, including the Karp--Lipton collapse under the stated assumptions; uniform polynomial-time construction would instead give a polynomial-time SAT algorithm.

Caching supplies an accounting bound for exact reuse. A terminating certified fallback procedure supplies correctness independently of predictor accuracy. A fixed policy yields expected savings when the solver work avoided on certified successes exceeds fast-path overhead and repays acquisition; explicit budgets bound additional work when prediction fails. These are elementary conditional guarantees, not evidence that an actual learned system satisfies their premises. A useful evaluation must jointly report acquisition, storage, inference, verification, fallback, and update costs against solver-only and cache-only baselines. The remaining practical question is under what conditions reusable knowledge repays the computation required to obtain it.

\FloatBarrier
\begin{thebibliography}{99}

\bibitem[Baker et~al.(1975)Baker, Gill, and Solovay]{baker1975}
Baker, T., Gill, J., and Solovay, R. (1975).
Relativizations of the $\mathcal P = ?\mathcal{NP}$ question.
\emph{SIAM Journal on Computing}, 4(4), 431--442.
\href{https://doi.org/10.1137/0204037}{doi:10.1137/0204037}.

\bibitem[Balc\'azar et~al.(1997)Balc\'azar, Gavald\`a, and Siegelmann]{balcazar1997neural}
Balc\'azar, J.~L., Gavald\`a, R., and Siegelmann, H.~T. (1997).
Computational power of neural networks: A characterization in terms of Kolmogorov complexity.
\emph{IEEE Transactions on Information Theory}, 43(4), 1175--1183.
\href{https://doi.org/10.1109/18.605580}{doi:10.1109/18.605580}.

\bibitem[Blum and Rivest(1988)]{blum1988}
Blum, A., and Rivest, R.~L. (1988).
Training a 3-node neural network is NP-complete.
\emph{Advances in Neural Information Processing Systems}, 1.
\url{https://proceedings.neurips.cc/paper/1988/hash/3def184ad8f4755ff269862ea77393dd-Abstract.html}.

\bibitem[Cadoli et~al.(2002)Cadoli, Donini, Liberatore, and Schaerf]{cadoli2002}
Cadoli, M., Donini, F.~M., Liberatore, P., and Schaerf, M. (2002).
Preprocessing of intractable problems.
\emph{Information and Computation}, 176(2), 89--120.
\href{https://doi.org/10.1006/inco.2001.3043}{doi:10.1006/inco.2001.3043}.

\bibitem[Cook(1971)]{cook1971}
Cook, S.~A. (1971).
The complexity of theorem-proving procedures.
In \emph{Proceedings of the Third Annual ACM Symposium on Theory of Computing}, pp.~151--158.
\href{https://doi.org/10.1145/800157.805047}{doi:10.1145/800157.805047}.

\bibitem[Cook and Reckhow(1979)]{cook1979}
Cook, S.~A., and Reckhow, R.~A. (1979).
The relative efficiency of propositional proof systems.
\emph{Journal of Symbolic Logic}, 44(1), 36--50.
\href{https://doi.org/10.2307/2273702}{doi:10.2307/2273702}.

\bibitem[Darwiche and Marquis(2002)]{darwiche2002}
Darwiche, A., and Marquis, P. (2002).
A knowledge compilation map.
\emph{Journal of Artificial Intelligence Research}, 17, 229--264.
\href{https://doi.org/10.1613/jair.989}{doi:10.1613/jair.989}.

\bibitem[Deolalikar(2010)]{deolalikar2010claim}
Deolalikar, V. (2010).
$P$ is not equal to $NP$.
Circulated manuscript, 6 August; failed separation claim.
\url{https://wscor.win.tue.nl/woeginger/P-versus-NP/Deolalikar.pdf}.

\bibitem[Fortnow(2022)]{fortnow2022fifty}
Fortnow, L. (2022).
Fifty years of P vs. NP and the possibility of the impossible.
\emph{Communications of the ACM}, 65(1), 76--85.
\href{https://doi.org/10.1145/3460351}{doi:10.1145/3460351}.

\bibitem[Fortnow(2026)]{fortnow2026manage}
Fortnow, L. (2026).
How does machine learning manage complexity?
Preprint, \emph{arXiv:2604.07233}; also ECCC TR26-053.
\url{https://arxiv.org/abs/2604.07233}.

\bibitem[Gershman and Goodman(2014)]{gershman2014}
Gershman, S.~J., and Goodman, N.~D. (2014).
Amortized inference in probabilistic reasoning.
In \emph{Proceedings of the 36th Annual Conference of the Cognitive Science Society}.
\url{https://escholarship.org/uc/item/34j1h7k5}.

\bibitem[Hertrich and Skutella(2023)]{hertrich2023knapsack}
Hertrich, C., and Skutella, M. (2023).
Provably good solutions to the knapsack problem via neural networks of bounded size.
\emph{INFORMS Journal on Computing}, 35(5), 1079--1097.
Earlier conference version at AAAI 2021.
\href{https://doi.org/10.1287/ijoc.2021.0225}{doi:10.1287/ijoc.2021.0225}.

\bibitem[Immerman(2010)]{immerman2010objections}
Immerman, N. (2010).
Technical objections to Deolalikar's claimed proof.
Letter published by R.~Lipton in ``Fatal Flaws in Deolalikar's Proof?'', 12 August.
\url{https://rjlipton.com/2010/08/12/fatal-flaws-in-deolalikars-proof/}.

\bibitem[Juba(2013)]{juba2013implicit}
Juba, B. (2013).
Implicit learning of common sense for reasoning.
In \emph{Proceedings of IJCAI}, pp.~939--945.
\url{https://ijcai.org/proceedings/2013/144}.

\bibitem[Kahneman and Klein(2009)]{kahneman2009}
Kahneman, D., and Klein, G. (2009).
Conditions for intuitive expertise: A failure to disagree.
\emph{American Psychologist}, 64(6), 515--526.
\href{https://doi.org/10.1037/a0016755}{doi:10.1037/a0016755}.

\bibitem[Karp and Lipton(1980)]{karp1980}
Karp, R.~M., and Lipton, R.~J. (1980).
Some connections between nonuniform and uniform complexity classes.
In \emph{Proceedings of the Twelfth Annual ACM Symposium on Theory of Computing}, pp.~302--309.
\href{https://doi.org/10.1145/800141.804678}{doi:10.1145/800141.804678}.

\bibitem[Karp and Lipton(1982)]{karp1982}
Karp, R.~M., and Lipton, R.~J. (1982).
Turing machines that take advice.
\emph{L'Enseignement Math\'ematique}, 28, 191--209.
\href{https://doi.org/10.5169/seals-52237}{doi:10.5169/seals-52237}.

\bibitem[Khardon and Roth(1997)]{khardon1997}
Khardon, R., and Roth, D. (1997).
Learning to reason.
\emph{Journal of the ACM}, 44(5), 697--725.
\href{https://doi.org/10.1145/265910.265918}{doi:10.1145/265910.265918}.

\bibitem[Lin et~al.(2021)Lin, Jaech, Li, Gormley, and Eisner]{lin2021limitations}
Lin, C.-C., Jaech, A., Li, X., Gormley, M.~R., and Eisner, J. (2021).
Limitations of autoregressive models and their alternatives.
In \emph{Proceedings of NAACL-HLT}, pp.~5147--5173.
\href{https://doi.org/10.18653/v1/2021.naacl-main.405}{doi:10.18653/v1/2021.naacl-main.405}.

\bibitem[McConnell et~al.(2011)McConnell, Mehlhorn, N\"aher, and Schweitzer]{mcconnell2011}
McConnell, R.~M., Mehlhorn, K., N\"aher, S., and Schweitzer, P. (2011).
Certifying algorithms.
\emph{Computer Science Review}, 5(2), 119--161.
\href{https://doi.org/10.1016/j.cosrev.2010.09.009}{doi:10.1016/j.cosrev.2010.09.009}.

\bibitem[Merrill and Sabharwal(2024)]{merrill2024cot}
Merrill, W., and Sabharwal, A. (2024).
The expressive power of transformers with chain of thought.
\emph{International Conference on Learning Representations}.
\url{https://arxiv.org/abs/2310.07923}.

\bibitem[Michie(1968)]{michie1968}
Michie, D. (1968).
``Memo'' functions and machine learning.
\emph{Nature}, 218, 19--22.
\href{https://doi.org/10.1038/218019a0}{doi:10.1038/218019a0}.

\bibitem[Minton(1988)]{minton1988utility}
Minton, S. (1988).
Quantitative results concerning the utility of explanation-based learning.
In \emph{Proceedings of AAAI}.
\url{https://cdn.aaai.org/AAAI/1988/AAAI88-100.pdf}.

\bibitem[Moses and Tennenholtz(1996)]{moses1996knowledge}
Moses, Y., and Tennenholtz, M. (1996).
Off-line reasoning for on-line efficiency: Knowledge bases.
\emph{Artificial Intelligence}, 83(2), 229--239.
\href{https://doi.org/10.1016/0004-3702(95)00015-1}{doi:10.1016/0004-3702(95)00015-1}.

\bibitem[Rigo(2014)]{rigo2014}
Rigo, M. (2014).
\emph{Formal Languages, Automata and Numeration Systems 2: Applications to Recognizability and Decidability}.
ISTE and Wiley. Chapter~1, Remark~1.12, p.~12.
\href{https://doi.org/10.1002/9781119042853}{doi:10.1002/9781119042853}.

\bibitem[Selsam and Bj\o rner(2019)]{selsam2019core}
Selsam, D., and Bj\o rner, N. (2019).
Guiding high-performance SAT solvers with unsat-core predictions.
In \emph{Theory and Applications of Satisfiability Testing (SAT 2019)}.
\url{https://arxiv.org/abs/1903.04671}.

\bibitem[Selsam et~al.(2019)Selsam, Lamm, B\"unz, Liang, de~Moura, and Dill]{selsam2019}
Selsam, D., Lamm, M., B\"unz, B., Liang, P., de~Moura, L., and Dill, D.~L. (2019).
Learning a SAT solver from single-bit supervision.
\emph{arXiv:1802.03685}, version~4 (first submitted 2018).
\url{https://arxiv.org/abs/1802.03685}.

\bibitem[Siegelmann and Sontag(1995)]{siegelmann1995}
Siegelmann, H.~T., and Sontag, E.~D. (1995).
On the computational power of neural nets.
\emph{Journal of Computer and System Sciences}, 50(1), 132--150.
\href{https://doi.org/10.1006/jcss.1995.1013}{doi:10.1006/jcss.1995.1013}.

\bibitem[Stuhlm\"uller et~al.(2013)Stuhlm\"uller, Taylor, and Goodman]{stuhlmuller2013}
Stuhlm\"uller, A., Taylor, J., and Goodman, N.~D. (2013).
Learning stochastic inverses.
\emph{Advances in Neural Information Processing Systems}, 26.
\url{https://proceedings.neurips.cc/paper/2013/hash/7f53f8c6c730af6aeb52e66eb74d8507-Abstract.html}.

\bibitem[Tadepalli and Natarajan(1996)]{tadepalli1996speedup}
Tadepalli, P., and Natarajan, B.~K. (1996).
A formal framework for speedup learning from problems and solutions.
\emph{Journal of Artificial Intelligence Research}, 4, 445--475.
\href{https://doi.org/10.1613/jair.154}{doi:10.1613/jair.154}.

\bibitem[Wen(2010a)]{wen2010mirrored}
Wen, H.~X. (2010a).
Mirrored language structure and innate logic of the human brain as a computable model of the oracle Turing machine.
\emph{arXiv:1006.2495}. Unestablished $P=NP$ claim.
\url{https://arxiv.org/abs/1006.2495}.

\bibitem[Wen(2010b)]{wen2010recognition}
Wen, H.~X. (2010b).
Knowledge recognition algorithm enables $P=NP$.
\emph{arXiv:1009.0884}. Unestablished $P=NP$ claim.
\url{https://arxiv.org/abs/1009.0884}.

\bibitem[Xu et~al.(2008)Xu, Hutter, Hoos, and Leyton-Brown]{xu2008satzilla}
Xu, L., Hutter, F., Hoos, H.~H., and Leyton-Brown, K. (2008).
SATzilla: Portfolio-based algorithm selection for SAT.
\emph{Journal of Artificial Intelligence Research}, 32, 565--606.
\href{https://doi.org/10.1613/jair.2490}{doi:10.1613/jair.2490}.

\bibitem[Yu et~al.(2024)Yu, Xu, Weston, and Kulikov]{yu2024}
Yu, P., Xu, J., Weston, J., and Kulikov, I. (2024).
Distilling System 2 into System 1.
\emph{arXiv:2407.06023}.
\url{https://arxiv.org/abs/2407.06023}.

\bibitem[Yun et~al.(2019)Yun, Sra, and Jadbabaie]{yun2019}
Yun, C., Sra, S., and Jadbabaie, A. (2019).
Small ReLU networks are powerful memorizers: A tight analysis of memorization capacity.
\emph{Advances in Neural Information Processing Systems}, 32.
\url{https://arxiv.org/abs/1810.07770}.

\end{thebibliography}
\end{document}
